\begin{topic}[id=profinet,
             difficulty={Mittel},
             type={Protokoll + Systemanalyse},
             prerequisites={Grundlagen industrieller Kommunikation; Basiswissen Ethernet hilfreich},
             practice={optional},
             owner={Dozententeam},
             updated={2026-04-09},
             tags={ Profinet, Real-Time, IRT, GSD, Diagnose, Security, Engineering }]{Profinet – Konzepte \& Einsatzgebiete}

\TopicDescription{Profinet verbindet Ethernet mit industriellen Anforderungen. Du lernst die Betriebsarten (RT/IRT), Engineering über GSD-Dateien, Diagnosepfade und Security-Aspekte (Netzsegmentierung, Zertifikate, Benutzerrechte). Wir betrachten typische Einsatzgebiete, Interoperabilität mit bestehenden Systemen und Trends (TSN, OPC UA-Kopplung). Ziel ist die Fähigkeit, Profinet-Topologien zu lesen, Risiken abzuschätzen und sinnvolle Integrationsmuster zu beschreiben.}

\TopicLeitfragen{\item Wann brauche ich IRT statt RT?
\item Wie sieht ein sinnvolles Segmentierungs-/Zertifikatskonzept aus?
\item Wie koppelt man Profinet sinnvoll an OPC UA?
}
\TopicScope{\item Keine vollständige Inbetriebnahme-Anleitung.
\item Kein Vendorvergleich.
}
\TopicPractice{Optionale Praxisidee: Beispiel-Topologie mit Diagnosewegen und Kopplungsskizze.}

\TopicKeywords{Profinet, Real-Time, IRT, GSD, Diagnose, Security, Engineering}

\nocite{piProfinet,siemensProfinet}
\end{topic}
